\documentclass[10pt,a4paper]{amsart}
\usepackage[margin=19mm]{geometry}
\usepackage{amsmath,amssymb,mathtools}
\usepackage[scr=rsfs,cal=euler]{mathalfa}
\usepackage{xfrac}
\usepackage[unicode,hidelinks,bookmarks=false]{hyperref}
\newcommand{\ie}{i.e.\ }
\newcommand{\cf}{cf.\ }
\newcommand{\resp}{respectively\ }
\newcommand{\Subsection}{\subsection}
\providecommand{\nobreakdash}{\nobreak\hbox{-}}
\setlength{\emergencystretch}{2em}
\multlinegap0pt
\newtheorem{assumption}{Assumption}
\newtheorem{proposition}{Proposition}
\newtheorem{lemma}{Lemma}
\def\R{\ensuremath{\mathbb{R}}}
\title{Natural Invariant Measures for Chaotic Game Dynamics: Finding Order in Chaos}
\author{Jakub Bielawski, Thiparat Chotibut, Fryderyk Falniowski, Micha{\l} Misiurewicz, Georgios Piliouras}
\date{Source: arXiv:2607.21805v1 (CC BY 4.0)}
\begin{document}
\maketitle

\noindent\emph{Source and licence.} Natural Invariant Measures for Chaotic Game Dynamics: Finding Order in Chaos. Version arXiv:2607.21805v1 (CC BY 4.0), distributed by arXiv under the Creative Commons Attribution 4.0 International licence, http://creativecommons.org/licenses/by/4.0/, as recorded in the arXiv licence metadata for this exact version and shown on the abstract page https://arxiv.org/abs/2607.21805v1. Full licence notice: \texttt{licenses/arxiv-2607.21805-CC-BY-4.0.txt}.\par\smallskip

\noindent\textit{Editorial scope.} Counted unit: the article's proposition averages2 (source0.tex lines 501--509). The article's complete original proof of it is delivered with it (source0.tex lines 511--550, one \texttt{proof} environment of 40 lines). No mathematics is written, repaired or re-derived: the statement and the proof are the article's own lines, in the article's own notation, and no retained line is substituted or re-flowed.  Retained uncounted as the source-local context the proof needs: lines 308--310; lines 568--572.  Nothing was omitted from any retained span, so the package carries no omission record.  No retained line contains a citation, so the wrapper carries no bibliography and no bibliography entry.  No retained line contains a figure environment, an image-inclusion command or drawing code, so no image is delivered and the package declares no asset dependency.  Editorial formatting only, in addition: an AMS \texttt{amsart} preamble that restates the article's own declarations and macros used by the retained lines, namely the environments \texttt{assumption}, \texttt{proposition}, \texttt{lemma} on separate counters, together with the standard packages those lines need.  The retained environments print in this document as Assumption 1, Proposition 1, Lemma 1 in order of appearance, whereas the article prints the same items with the numbering of its own full document.  The complete native source of the article as distributed in the arXiv e-print for this exact version is bundled unchanged as \texttt{sources/205893/source0.tex}.
\begin{assumption}\label{as1}
Let $(X,f)$ be a dynamical system, that is $X$ is nonempty and compact and $f\colon X \to X$ is a continuous map. Let $O\colon X\to\R$ be an observable (a continuous function).
\end{assumption}

\begin{lemma}\label{split}
If a sequence $(\gamma_i)_{i \in \mathbb{N}} \subset \R$ can be split
into subsequences such that all these subsequences converge to the
same point $z$, then $(\gamma_i)$ also converges to $z$.
\end{lemma}

\begin{proposition}\label{averages2}
%Let $X$ be a compact metric space, $f:X\to X$ a continuous map, $O:X\to\R$ a continuous function (an observable).
Let Assumption \ref{as1} hold. Assume that the dynamical system $(X,f)$ has an attracting periodic orbit of period $n > 0$ that attracts almost all initial states.
Then for almost all points $x \in X$
\[
\lim_{i\rightarrow \infty} \hat{O}(f^i(x)) = \hat{O}(z),
\]
where $z$ is an arbitrary point of the attracting periodic orbit.
\end{proposition}

\begin{proof}
We are going to show the proof in two steps.

Step 1. We will show that if $n$ divides $m$, then for almost all initializations, the limit of the $m$-th iterate time-averaged observations for any continuous function $O$ is captured by its time average over the attracting periodic orbit, or in other words, that
for almost all points $x \in X$
\begin{equation}\label{thm_av-eq1}
\lim_{i \rightarrow \infty} \hat{O}_m(f^i(x)) =  \hat{O}_m(z).
\end{equation}
Clearly, if $n$ divides $m$ then for every choice of the element $z$
of the periodic orbit the averages $\hat{O}_m(z)$ are equal to each
other. Let $x_i = f^i(x)$ be a trajectory attracted to the orbit of
$z$. It means that the sequence $(x_i)$ can be split into disjoint
subsequences $(x_{i_j})_{j \in \{1,2,\ldots,n\}}$ such that each
$(x_{i_j})$ converges to an element of the attracting periodic orbit
of $z$. Because the observable $O$ is continuous, we have that the
subsequences $\hat{O}_m(x_{i_j})$, for $j \in \{1,2,\ldots,n\}$,
converge to the same limit $\hat{O}_m(z)$. 
%As a sequence $(\gamma_i)_{i \in \mathbb{N}} \subset \R$ can be split into subsequences such that all these subsequences converge to the same point $z$, %then $(\gamma_i)$ also converges to $z$.
%Lemma \ref{split} we get that 
Thus, $\hat{O}_m(x_{i})$ also converges to $\hat{O}_m(z)$.

Step 2. Let $\varepsilon > 0$ and let $x_i = f^i(x)$ be a trajectory attracted to the orbit of
$z$. We are going to show that
\[
 | \hat{O}(z) - \hat{O}(x_i) | < \varepsilon
\]
for $i$ sufficiently large.

Choose $m>0$ such that $m=0 \mod (n)$.  Then because $z$ is an element of the periodic orbit of period $n$ we have that $\hat{O}(z) = \hat{O}_m(z)$.
Second, from the definition of the time-average $\hat{O}$ we get that
\begin{equation}\label{thm_av-eq2}
 | \hat{O}_m (x_i) - \hat{O} (x_i) | < \frac{\varepsilon}{2}
\end{equation}
for $m$ sufficiently large. Therefore, by combining \eqref{thm_av-eq1}, \eqref{thm_av-eq2} and the fact that $\hat{O}(z) = \hat{O}_m(z)$ we obtain
\begin{align*}
 \big| \hat{O} (z) - \hat{O} (x_i) \big| &= \big| \hat{O} (z) - \hat{O}_m (x_i) + \hat{O}_m (x_i) - \hat{O} (x_i) \big| \\
 &\leqslant \big| \hat{O}_m (z) - \hat{O}_m (x_i) \big| +  \big| \hat{O}_m (x_i) - \hat{O} (x_i) \big| < \frac{\varepsilon}{2} + \frac{\varepsilon}{2} = \varepsilon
\end{align*}
for $i$ and $m$ sufficiently large.
\end{proof}

\end{document}
